\documentclass[11pt]{article}
\usepackage[T1]{fontenc}
\usepackage{lmodern}
\usepackage[margin=1in]{geometry}
\usepackage{amsmath}
\usepackage{booktabs}
\usepackage{array}
\usepackage{microtype}
\usepackage[hidelinks]{hyperref}
% Generated by export_results.py. Do not edit numerical results by hand.
\newcommand{\StudyRun}{study-20261026T145746Z-ec1e8589}
\newcommand{\PrimarySourceHash}{48b70bf9170e3b3b367bb545215f9778ee30133b748fdc969feaabac32eec7e0}
\newcommand{\PrimaryUnits}{132}
\newcommand{\PrimaryWorkflows}{396}
\newcommand{\SpecificationCount}{30}
\newcommand{\VariantCount}{66}
\newcommand{\MatchedIntentCount}{17}
\newcommand{\PrimaryRequests}{535}
\newcommand{\PrimaryKnownRequests}{534}
\newcommand{\PrimaryUnknownRequests}{1}
\newcommand{\PrimaryInputTokens}{2,136,760}
\newcommand{\PrimaryOutputTokens}{882,428}
\newcommand{\PrimaryTotalTokens}{3,019,188}
\newcommand{\PrimaryOperationalWorkflows}{3}
\newcommand{\ModelOne}{gpt-5.3-codex}
\newcommand{\ModelTwo}{gpt-6.1-sol}
\newcommand{\IterationBudget}{5}
\newcommand{\OutputTokenBudget}{12,000}
\newcommand{\AWSInitialCompiled}{32}
\newcommand{\AWSInitialPassed}{0}
\newcommand{\AWSInitialUnits}{36}
\newcommand{\AWSFeedbackCompiled}{36}
\newcommand{\AWSFeedbackPassed}{30}
\newcommand{\AWSFeedbackUnits}{36}
\newcommand{\AWSInitialFlowFailures}{32}
\newcommand{\AWSInitialDefaultSGFailures}{32}
\newcommand{\PooledResultRows}{%
Azure/Bicep & S & 58 & 54/58 (93.1\%) & 50/58 (86.2\%) \\
Azure/Bicep & B & 58 & 55/58 (94.8\%) & 52/58 (89.7\%) \\
Azure/Bicep & F & 58 & 56/58 (96.6\%) & 56/58 (96.6\%) \\
Azure/Terraform & S & 38 & 37/38 (97.4\%) & 12/38 (31.6\%) \\
Azure/Terraform & B & 38 & 34/38 (89.5\%) & 18/38 (47.4\%) \\
Azure/Terraform & F & 38 & 38/38 (100.0\%) & 38/38 (100.0\%) \\
AWS/Terraform & S & 36 & 32/36 (88.9\%) & 0/36 (0.0\%) \\
AWS/Terraform & B & 36 & 34/36 (94.4\%) & 0/36 (0.0\%) \\
AWS/Terraform & F & 36 & 36/36 (100.0\%) & 30/36 (83.3\%) \\
}
\newcommand{\ModelResultRows}{%
Azure/Bicep & M1 & S & 29 & 28/29 (96.6\%) & 26/29 (89.7\%) \\
Azure/Bicep & M1 & B & 29 & 28/29 (96.6\%) & 27/29 (93.1\%) \\
Azure/Bicep & M1 & F & 29 & 28/29 (96.6\%) & 28/29 (96.6\%) \\
Azure/Bicep & M2 & S & 29 & 26/29 (89.7\%) & 24/29 (82.8\%) \\
Azure/Bicep & M2 & B & 29 & 27/29 (93.1\%) & 25/29 (86.2\%) \\
Azure/Bicep & M2 & F & 29 & 28/29 (96.6\%) & 28/29 (96.6\%) \\
Azure/Terraform & M1 & S & 19 & 19/19 (100.0\%) & 7/19 (36.8\%) \\
Azure/Terraform & M1 & B & 19 & 17/19 (89.5\%) & 12/19 (63.2\%) \\
Azure/Terraform & M1 & F & 19 & 19/19 (100.0\%) & 19/19 (100.0\%) \\
Azure/Terraform & M2 & S & 19 & 18/19 (94.7\%) & 5/19 (26.3\%) \\
Azure/Terraform & M2 & B & 19 & 17/19 (89.5\%) & 6/19 (31.6\%) \\
Azure/Terraform & M2 & F & 19 & 19/19 (100.0\%) & 19/19 (100.0\%) \\
AWS/Terraform & M1 & S & 18 & 16/18 (88.9\%) & 0/18 (0.0\%) \\
AWS/Terraform & M1 & B & 18 & 16/18 (88.9\%) & 0/18 (0.0\%) \\
AWS/Terraform & M1 & F & 18 & 18/18 (100.0\%) & 13/18 (72.2\%) \\
AWS/Terraform & M2 & S & 18 & 16/18 (88.9\%) & 0/18 (0.0\%) \\
AWS/Terraform & M2 & B & 18 & 18/18 (100.0\%) & 0/18 (0.0\%) \\
AWS/Terraform & M2 & F & 18 & 18/18 (100.0\%) & 17/18 (94.4\%) \\
}
\newcommand{\MatchedResultRows}{%
Azure/Bicep & S & 34 & 33/34 (97.1\%) & 31/34 (91.2\%) \\
Azure/Bicep & B & 34 & 34/34 (100.0\%) & 32/34 (94.1\%) \\
Azure/Bicep & F & 34 & 34/34 (100.0\%) & 34/34 (100.0\%) \\
Azure/Terraform & S & 34 & 33/34 (97.1\%) & 10/34 (29.4\%) \\
Azure/Terraform & B & 34 & 30/34 (88.2\%) & 15/34 (44.1\%) \\
Azure/Terraform & F & 34 & 34/34 (100.0\%) & 34/34 (100.0\%) \\
AWS/Terraform & S & 34 & 30/34 (88.2\%) & 0/34 (0.0\%) \\
AWS/Terraform & B & 34 & 32/34 (94.1\%) & 0/34 (0.0\%) \\
AWS/Terraform & F & 34 & 34/34 (100.0\%) & 28/34 (82.4\%) \\
}
\newcommand{\PairedResultRows}{%
Azure/Bicep & M1 & 29 & 1 & 1 & 0 & 27 & +3.4 \\
Azure/Bicep & M2 & 29 & 1 & 3 & 0 & 25 & +10.3 \\
Azure/Terraform & M1 & 19 & 0 & 7 & 0 & 12 & +36.8 \\
Azure/Terraform & M2 & 19 & 0 & 13 & 0 & 6 & +68.4 \\
AWS/Terraform & M1 & 18 & 5 & 13 & 0 & 0 & +72.2 \\
AWS/Terraform & M2 & 18 & 1 & 17 & 0 & 0 & +94.4 \\
}
\newcommand{\AWSPolicyRows}{%
M1 & S & 18 & 16 & 16 & 16 & 0 & 1.00 \\
M1 & B & 18 & 16 & 16 & 16 & 0 & 5.00 \\
M1 & F & 18 & 18 & 2 & 1 & 13 & 4.39 \\
M2 & S & 18 & 16 & 16 & 16 & 0 & 1.00 \\
M2 & B & 18 & 18 & 18 & 16 & 0 & 5.00 \\
M2 & F & 18 & 18 & 0 & 0 & 17 & 3.22 \\
}
\newcommand{\UniqueUsageRows}{%
M1 & 264 & 0 & 1,052,495 & 524,628 & 1,577,123 \\
M2 & 270 & 1 & 1,084,265 & 357,800 & 1,442,065 \\
}
\newcommand{\AttributedUsageRows}{%
M1 & S & 66 & 66 & 0 & 193,295 & 110,285 & 42.3 \\
M1 & B & 66 & 185 & 0 & 661,711 & 344,851 & 122.4 \\
M1 & F & 66 & 145 & 0 & 584,079 & 290,062 & 101.4 \\
M2 & S & 66 & 65 & 1 & 205,966 & 78,075 & 43.2 \\
M2 & B & 66 & 206 & 1 & 770,410 & 236,789 & 119.0 \\
M2 & F & 66 & 129 & 1 & 519,821 & 199,086 & 90.3 \\
}
\newcommand{\RecoveredRequests}{542}
\newcommand{\RetryNewRequests}{7}
\newcommand{\RecoveredOperationalWorkflows}{0}
\newcommand{\RecoveredTotalTokens}{3,158,910}
\newcommand{\RecoveredResultRows}{%
Azure/Bicep & S & 58 & 55/58 (94.8\%) & 50/58 (86.2\%) \\
Azure/Bicep & B & 58 & 56/58 (96.6\%) & 52/58 (89.7\%) \\
Azure/Bicep & F & 58 & 57/58 (98.3\%) & 57/58 (98.3\%) \\
Azure/Terraform & S & 38 & 37/38 (97.4\%) & 12/38 (31.6\%) \\
Azure/Terraform & B & 38 & 34/38 (89.5\%) & 18/38 (47.4\%) \\
Azure/Terraform & F & 38 & 38/38 (100.0\%) & 38/38 (100.0\%) \\
AWS/Terraform & S & 36 & 32/36 (88.9\%) & 0/36 (0.0\%) \\
AWS/Terraform & B & 36 & 34/36 (94.4\%) & 0/36 (0.0\%) \\
AWS/Terraform & F & 36 & 36/36 (100.0\%) & 30/36 (83.3\%) \\
}


\title{From Human-Written Network Requirements to Infrastructure as Code:\\
An Exploratory Paired Study of Validator Feedback Across Azure and AWS}
\author{Author information to be completed}
\date{Working draft --- \today}

\begin{document}
\maketitle
\begin{center}
\small\textbf{Status: exploratory, artifact-backed working draft; not submission-ready.}\\
Static validator outcomes are measured. Requirement satisfaction is not yet assessed.
\end{center}

\begin{abstract}
Generating infrastructure as code (IaC) from human-written requirements involves
more than producing compilable configuration. We examine this distinction in an
exploratory study of \SpecificationCount{} network specifications expanded into
\VariantCount{} supported Azure/Bicep, Azure/Terraform, and AWS/Terraform variants.
Two configured model identifiers produce \PrimaryUnits{} paired experimental units
and \PrimaryWorkflows{} condition workflows. Each unit shares an initial candidate
across one-shot generation, blind refinement, and refinement with compiler and
static-policy diagnostics. Refinement is bounded by \IterationBudget{} attempts.
The primary analysis retains the original completed outcomes, including operational
failures, rather than replacing them after inspecting results.
AWS/Terraform illustrates a large gap between compilation and policy convergence:
\AWSInitialCompiled{}/\AWSInitialUnits{} initial candidates compile, but
\AWSInitialPassed{}/\AWSInitialUnits{} pass all enabled visible validators.
Feedback-guided workflows reach \AWSFeedbackPassed{}/\AWSFeedbackUnits{}.
All compiling initial AWS candidates fail the same two policies concerning flow
logging and default-security-group restrictions. This pattern occurs in both models
and is consistent with policy requirements absent from some written scenarios;
it does not establish a model-specific cause. We report paired outcomes,
model-stratified results, token consumption, and an audited operational-retry
sensitivity analysis. The observations support bounded claims about convergence
to the configured validators, not semantic requirement satisfaction, deployment
correctness, general cloud difficulty, or model superiority.
\end{abstract}

\section{Introduction}
IaC generation from requirements can fail at several distinct levels: language
validity, provider-schema consistency, policy compliance, interpretation of the
requested design, deployability, and runtime behavior. Treating these levels as a
single acceptance score can obscure what an experiment actually measures.
Terraform's validation command, for example, checks syntax and internal consistency
but does not validate remote services or provider APIs~\cite{terraformvalidate}.
Likewise, Checkov checks a collection of static misconfiguration policies rather
than proving end-to-end network behavior~\cite{checkov}.

Our motivating goal is to understand where models struggle when translating
human-readable network requirements. The present experiment addresses only an
initial, measurable part of that goal: whether visible validation diagnostics
help bounded refinement converge to the configured static checks. Independent
assessment of the requirements themselves remains future work.

The contributions of this working draft are:
\begin{enumerate}
  \item A file-backed, requirements-only network corpus and a paired workflow
  design that separates initial generation, blind refinement, and diagnostic
  feedback.
  \item An artifact-backed descriptive analysis across three cloud/IaC targets,
  with explicit model stratification, denominators, and shared-candidate accounting.
  \item A policy-level examination of the AWS/Terraform gap and a clear separation
  between validator convergence, operational reliability, and unassessed semantics.
\end{enumerate}
We do not claim to invent feedback-driven IaC generation or to establish
confirmatory statistical findings.

\section{Related Work and Measurement Scope}
Palavalli and Santolucito investigate feedback loops for LLM-based IaC
generation~\cite{palavalli2024feedback}. Their work establishes that returning
errors and warnings to a generator is prior art. Our potential contribution must
therefore rest on the paired controls, network-requirements corpus, multi-target
analysis, and the distinction between policy convergence and requirement
satisfaction, not on the existence of a repair loop.

The literature review is preliminary and is not a systematic review. A
venue-specific related-work expansion is necessary before submission.
The official Terraform and Checkov descriptions~\cite{terraformvalidate,checkov}
also motivate an important measurement boundary: static validation success is
neither a deployment test nor a guarantee that a generated system realizes a
user's intended architecture or connectivity.

\section{Research Questions}
\paragraph{RQ1: Static convergence.}
Within the saved scenario--model pairs, how do final compilation and
schema-validation results and
visible-validator outcomes differ between diagnostic feedback and blind
refinement under the same maximum attempt budget?
\paragraph{RQ2: Target and model patterns.}
What descriptive differences appear after stratifying by cloud/IaC target and
configured model identifier, and what policy failures explain the observed
AWS/Terraform pattern?
\paragraph{RQ3: Resource use.}
What observed token and workflow-latency costs accompany these outcomes,
accounting for shared initial candidates and unknown usage?

These are exploratory questions. The AWS policy analysis was prompted by the
observed results and is explicitly post hoc. No preregistration, confirmatory
hypothesis test, population-level effect, or causal model comparison is claimed.

\section{Study Design}
\subsection{Corpus and targets}
The corpus consists of \SpecificationCount{} authored YAML specifications, not
requirements sampled from users or organizations. Specifications describe
objectives and acceptance, connectivity, security, and architectural requirements.
Their specificity varies: some prescribe components, addressing, or ports.
Calling them human-written does not make this a human-subjects study or establish
that they represent the distribution of real organizational requests.

The target expansion yields 29 Azure/Bicep, 19 Azure/Terraform, and 18
AWS/Terraform variants. Original scenarios NET-01--NET-10 retain their
Azure/Bicep scope. Newer scenarios declare supported targets in their source files.
Unsupported pairs are exclusions, not failures. Consequently, raw cross-target
rates mix different scenario sets. An appendix reports the
\MatchedIntentCount{} intents shared by all three targets, but even that subset
does not make the target-specific policy suites equivalent.

Prompts render the resolved specifications as prose and requirement lists. They
do not append internal resource graphs, exact-resource answer keys, or JSON/YAML
solution examples. The model still receives target-specific output constraints.
An earlier resource-graph expansion was withdrawn because it changed the task
into low-level graph translation. That expansion and earlier contract-fed pilots
are excluded from this paper's primary dataset.

\subsection{Experimental unit and conditions}
The unit is scenario variant $\times$ model identifier $\times$ repetition.
Two identifiers, \texttt{\ModelOne{}} (M1) and \texttt{\ModelTwo{}} (M2), each have
one repetition per supported variant. These are service identifiers, not
independently verified architecture or immutable model-weight versions.

Each unit generates one common initial candidate, then evaluates the paired
conditions in Table~\ref{tab:conditions}. The common candidate is not three
independent model generations. Subsequent refinement sessions are fresh,
tool-free Copilot SDK sessions receiving the scenario and previous candidate;
only the feedback condition receives visible diagnostic reports.

\begin{table}[tbp]
\centering\small
\begin{tabular}{>{\raggedright\arraybackslash}p{0.20\linewidth}>{\raggedright\arraybackslash}p{0.17\linewidth}>{\raggedright\arraybackslash}p{0.52\linewidth}}
\toprule
Condition & Maximum attempts & Revision information \\
\midrule
One-shot (S) & 1 & Initial scenario and target/output constraints. \\
Blind refinement (B) & \IterationBudget{} & Scenario and previous candidate, without validator diagnostics. \\
Tool feedback (F) & \IterationBudget{} & Same inputs plus compiler/schema and enabled policy diagnostics. \\
\bottomrule
\end{tabular}
\caption{Conditions. All three share the same initial candidate within a unit.
The refinement arms have equal maximum, not equal realized, budgets.}
\label{tab:conditions}
\end{table}

Each arm stops on a visible-validator pass, generation failure, or exhaustion of
the attempt budget. Blind refinement receives no diagnostic content, although
the decision to request another revision implicitly reveals non-convergence.
The visible checks are therefore both an intervention input and a stopping
criterion. We measure the final candidate under that rule, not an oracle-selected
best candidate. The one-shot comparison additionally changes the available budget;
the feedback-versus-blind comparison is the closer diagnostic-information control.

Scenario/model jobs and arm order are shuffled with scheduler seed 1729.
This seed does not control service-side sampling. Temperature and top-$p$ are
unset. The configured output limit is \OutputTokenBudget{} tokens per generation.
Copilot SDK 1.0.22 is used with workspace context, tools, memory, and local
instruction discovery disabled. Provider foundation instructions and service
policies remain outside experimental control.

\subsection{Validation and implementation constraints}
Bicep CLI 0.48.1 compiles Azure/Bicep to ARM. Terraform 1.9.8 validates against
the pinned AzureRM 4.0.1 and AWS 5.70.0 providers, with backend use disabled for
initialization. No Terraform plan/apply or cloud deployment is performed.
Checkov 3.3.29 scans the candidate when compilation/schema validation succeeds.
Conftest/OPA is disabled for this mixed-target run.

The generation harness imposes a restricted Terraform output language: resource
blocks, literal values, and direct references; it supplies the pinned provider
prelude. Modules, data sources, variable/local/output blocks, executable
provisioners, and dynamic expansion are not supported. These restrictions limit
external validity and can prevent otherwise reasonable Terraform designs.

The visible-pass outcome requires every enabled visible validator to pass;
disabled validators are neutral. A Checkov pass in this run requires a nonzero
passed-check count, no reported failed checks, and no parsing errors. It is not
a complete proof of security. Policy applicability and coverage depend on both
the target and the resources the model chooses to emit.

The explicit assessment mode leaves both requirement acceptance and independent
structural acceptance \emph{unassessed}. No exact-resource oracle grades even the
original NET scenarios in this run. No deployment, reachability, DNS-resolution,
cost-of-infrastructure, or requirement-completeness result is inferred.

\subsection{Evidence, interruptions, and analysis populations}
The evidence consists of saved prompts, responses, candidate files, validator
reports, per-generation usage metadata, complete paired results, and managed
checkpoints. The run identifier is:
\begin{center}\small\texttt{\StudyRun{}}\end{center}
A Windows checkpoint-replacement interruption required recovery during collection.
Fifty-four checkpointed units and one fully saved unit were preserved before
resuming. The completed-unit recovery was not a model regeneration.

The primary analysis uses the original completed checkpoint before an explicitly
authorized operational retry. It includes one failed shared generation in
NET-10/M2, affecting \PrimaryOperationalWorkflows{} condition workflows.
These remain operational failures in the original denominators, not evidence of
incorrect requirement interpretation.

The later retry is a separate sensitivity analysis. It repeats only that failed
paired unit, with its original settings and arm order, preserves the other 131
paired results, and retains the original evidence and failed request in accounting.
It is not a new independent repetition. We report original and replacement-based
results separately rather than silently promoting a repaired dataset to the
primary analysis.

\subsection{Metrics and descriptive analysis}
For target/model stratum $g$ and condition $a$, let $n_g$ be the number of saved
paired units, including operational failures. We report
\[
 \widehat p_{g,a}^{\,V}
 = \frac{\sum_{i\in g}\mathbf{1}(\text{final visible pass}_{i,a})}{n_g},
 \qquad
 \widehat p_{g,a}^{\,C}
 = \frac{\sum_{i\in g}\mathbf{1}(\text{final compilation/schema pass}_{i,a})}{n_g}.
\]
Missing compilation evidence from a generation failure is unknown, not an
observed compiler rejection; the pass-rate denominator nevertheless retains
that attempted unit. Acceptance has no percentage because it has not been assessed.

For feedback versus blind refinement, $N_{01}$ denotes pairs where only feedback
passes and $N_{10}$ pairs where only blind refinement passes. The descriptive
paired difference is
\[
 \widehat\Delta_g = 100\,\frac{N_{01}-N_{10}}{n_g}
 \quad\text{percentage points}.
\]
Pairs are not independent across related intents, targets, or specification
families. One service generation per initial scenario/model cell does not estimate
sampling variability. We therefore do not attach $p$-values, independence-based
confidence intervals, or significance labels to these exploratory counts.

\section{Results}
\subsection{Compilation and visible-validator convergence}
Table~\ref{tab:pooled} pools the two models by summing counts and denominators,
not by averaging percentages. Model-stratified results are given in
Appendix~\ref{app:models}; they are necessary before discussing model differences.
Each target has a different scenario composition and policy surface.

\begin{table}[tbp]
\centering\small
\begin{tabular}{llrrr}
\toprule
Target & Arm & $n$ & Compilation/schema pass & Visible pass \\
\midrule
\PooledResultRows
\bottomrule
\end{tabular}
\caption{Original completed outcomes, pooled across M1 and M2.
S: one-shot; B: blind; F: feedback. The original shared generation failure has
unknown compilation status in all three workflows and remains in the denominator.}
\label{tab:pooled}
\end{table}

The AWS/Terraform initial candidates demonstrate that compilation alone can
produce a misleading impression of success:
\AWSInitialCompiled{}/\AWSInitialUnits{} compile but none passes every enabled
visible check. Feedback reaches \AWSFeedbackPassed{}/\AWSFeedbackUnits{} visible
passes. This is a result about the configured validator workflow, not evidence
that those artifacts satisfy all stated requirements or work after deployment.

\begin{table}[tbp]
\centering\small
\begin{tabular}{llrrrrrr}
\toprule
Target & Model & $n$ & $N_{00}$ & $N_{01}$ & $N_{10}$ & $N_{11}$ & $\widehat\Delta$ (pp) \\
\midrule
\PairedResultRows
\bottomrule
\end{tabular}
\caption{Paired visible outcomes: blind refinement first, feedback second.
$00$: neither passes; $01$: feedback only; $10$: blind only; $11$: both pass.
These are descriptive within-stratum contrasts, not significance tests.}
\label{tab:paired}
\end{table}

\subsection{What explains the AWS pattern?}
Table~\ref{tab:aws} counts \emph{workflows} with at least one failure of each
policy, deduplicating repeated resource-level occurrences of the same check.
Every compiling initial AWS candidate fails
\texttt{CKV2\_AWS\_11} (VPC flow logging) and
\texttt{CKV2\_AWS\_12} (default-security-group restrictions):
\AWSInitialFlowFailures{} and \AWSInitialDefaultSGFailures{} workflows,
respectively. Compiler failures are not scanned and must not be interpreted as
having passed these policies.

\begin{table}[tbp]
\centering\small
\setlength{\tabcolsep}{5pt}
\begin{tabular}{llrrrrrr}
\toprule
Model & Arm & $n$ & Compiled & Flow-log fail & Default-SG fail & Visible pass & Mean attempts \\
\midrule
\AWSPolicyRows
\bottomrule
\end{tabular}
\caption{AWS/Terraform policy audit by model and condition.
Failure counts indicate workflows containing the named violation, not independent
policy tests. Unlisted policies can also cause a visible failure.}
\label{tab:aws}
\end{table}

This common baseline across both identifiers does not isolate a model effect.
The feedback outcomes differ descriptively between M1 and M2, but each has only
18 initial AWS generations and one repetition per scenario. The comparison is
also affected by stochastic generation, service timing, and policy applicability.

A concrete measurement concern is the difference between written requirements
and implicit policy expectations. NET-11 asks for a private isolated workload
network without public addressing; it does not explicitly require flow-log
infrastructure or management of the VPC's default security group. The policy
suite imposes these expectations nonetheless. Such policies may be reasonable
organizational baselines, but here they are not uniformly part of the initial
specification. Feedback can therefore teach the evaluator's implicit requirements,
not merely correct the model's interpretation of the supplied requirements.

This is evidence consistent with an evaluator-alignment explanation, not a
causal identification of the root cause. Establishing causality would require,
for example, controlled specification/policy-disclosure ablations and independent
review of equivalent valid designs. It would be premature to conclude that AWS
is intrinsically harder or that a model lacks AWS knowledge.

\subsection{Tokens and recorded latency}
The original completed run contains \PrimaryRequests{} saved generation requests:
\PrimaryKnownRequests{} have reported usage and \PrimaryUnknownRequests{} has
unknown usage. The observed subtotal is \PrimaryInputTokens{} input plus
\PrimaryOutputTokens{} output tokens (\PrimaryTotalTokens{} total).
Table~\ref{tab:unique} counts each saved request once.

\begin{table}[tbp]
\centering\small
\begin{tabular}{lrrrrr}
\toprule
Model & Known requests & Unknown & Input tokens & Output tokens & Total tokens \\
\midrule
\UniqueUsageRows
\bottomrule
\end{tabular}
\caption{Primary unique-request accounting across all targets, including the
operational failure. Unknown usage is not imputed as zero.}
\label{tab:unique}
\end{table}

\begin{table}[tbp]
\centering\small
\setlength{\tabcolsep}{4pt}
\begin{tabular}{llrrrrrr}
\toprule
Model & Arm & Workflows & Known req. & Unknown & Input tokens & Output tokens & Mean seconds \\
\midrule
\AttributedUsageRows
\bottomrule
\end{tabular}
\caption{Workflow-attributed usage and latency across the mixed-target workload.
Each arm includes the shared initial generation. \emph{Do not sum these rows as
an execution or billing total.} These pooled accounting figures are not
target-controlled performance comparisons.}
\label{tab:attributed}
\end{table}

Recorded workflow latency attributes the shared initial work to each condition
and includes generation, SDK overhead, and validation; it is not pure model
inference time. Whole-run wall time additionally includes operational interruption
and idle time, so it is not used as a model-speed measure.
Copilot allowance/credit accounting is not an Azure token-price invoice.
No monetary cost is inferred. SDK/backend retries may involve additional
unobserved physical inference calls beyond the saved request ledger.

\subsection{Operational-retry sensitivity}
The authorized retry uses \RetryNewRequests{} additional generation requests.
The replacement-based dataset has \RecoveredRequests{} saved requests and
\RecoveredOperationalWorkflows{} operational-failure workflows. Unique accounting
still includes the original failed request and has an observed token subtotal of
\RecoveredTotalTokens{}; the previously unknown usage remains unknown.
Table~\ref{tab:retry} reports replacement-based outcomes separately from
Table~\ref{tab:pooled}. The replacement changes no other paired result and must
not be counted as a second independent NET-10/M2 observation.

\begin{table}[tbp]
\centering\small
\begin{tabular}{llrrr}
\toprule
Target & Arm & $n$ & Compilation/schema pass & Visible pass \\
\midrule
\RecoveredResultRows
\bottomrule
\end{tabular}
\caption{Sensitivity analysis after the audited operational replacement.
Original outcomes and failed-request accounting remain available in the audit.}
\label{tab:retry}
\end{table}

\section{Threats to Validity}
\paragraph{Construct validity.}
Visible pass is a narrow static-convergence metric. The same validators supply
feedback and determine the outcome; there is no held-out semantic assessment.
Policy compliance can add infrastructure not requested by a scenario, while
resource omission can alter which checks are applicable. Neither a static pass
nor the absence of reported violations proves requirement completeness or
security. The paper does not score human ambiguity, underspecification, or
alternative valid architectures.

\paragraph{Internal validity.}
Sharing the initial candidate supports paired comparisons but correlates all
three condition outcomes. Early stopping makes realized token and attempt
budgets unequal. Blinding only diagnostic content does not remove the implicit
non-convergence signal from a request for another revision. Fresh sessions and
shuffled ordering reduce some contamination and timing effects, but they do not
control provider sampling, model updates, or all time-of-day effects. Operational
recovery and the explicitly authorized post-run retry must remain auditable.

\paragraph{External validity.}
The corpus is authored and small, not a representative requirements sample.
Scenarios overlap conceptually and through inheritance. The target matrix is
unbalanced, resource-language restrictions constrain feasible designs, and
Checkov's AWS, Azure/Terraform, and ARM policy surfaces are not equivalent.
Two model identifiers do not represent all models or deployment environments.
No real application, cloud deployment, production workload, or human-user study
is evaluated.

\paragraph{Conclusion validity.}
There is one initial generation per scenario/model cell, no power analysis, and
no preregistered confirmatory analysis. Apparent differences between 18-case AWS
strata must not be called significant or generalized as model superiority.
Post-hoc discovery of recurring policies motivates hypotheses rather than
confirming them. Known-token subtotals and composite latency must not be
misrepresented as full inference usage, dollar cost, or provider-independent speed.

\section{What Would Make the Study Submission-Ready?}
A stronger next study should freeze its corpus, policy disclosure, primary
endpoint, and analysis before collecting a new dataset. We propose:
\begin{enumerate}
  \item An independent requirements rubric accepting semantically valid alternative
  designs, with blinded expert assessment, disagreement resolution, and reported
  inter-rater agreement. A high-level requirement must not secretly demand one
  exact resource name or topology.
  \item A controlled ambiguity/underspecification study with clearly authored
  variants of the same intent. This is needed to address limitations of
  human-written requirements directly, rather than inferring them from compiler
  or policy failures.
  \item A policy-disclosure ablation: requirements alone versus requirements plus
  an explicitly documented policy baseline, with held-out assessment in both cases.
  This would test the evaluator-alignment explanation of the AWS pattern.
  \item Matched cross-target intents, additional independent repetitions, and a
  justified precision or power plan. Analysis should account for paired arms,
  repeated intents, families, target variants, and multiple comparisons.
  \item Safe deployment/runtime checks on a reviewed subset, measuring reachability,
  DNS behavior, exposure, and other relevant semantics. These remain separate
  endpoints, not automatic consequences of a static pass.
  \item A reviewed artifact release with exact dependency versions, policy scope,
  hashes, raw operational failures, and retry provenance. Service identifiers
  and usage coverage must be retained without credentials.
\end{enumerate}
These are proposed validation steps, not experiments already performed.

\section{Conclusion}
The defensible current finding is a descriptive, paired association between
diagnostic feedback and convergence to the configured static validators in a
bounded network-IaC workload. Compilation and policy convergence differ sharply,
especially for AWS/Terraform. Recurring policy failures across both models make
an implicit-policy explanation plausible and undermine a simple attribution to
one model. Independent requirement satisfaction remains unassessed. The current
artifact is suitable as the basis for an exploratory workshop/tool-study draft
or a pilot section, subject to venue expectations and further review; it is not
yet a confirmatory evaluation of human-requirements interpretation.

\section*{Reproducibility, Availability, and Draft Disclosure}
The accompanying exporter derives the tables from saved checkpoints and final
validator reports without executing model requests or deploying infrastructure.
It records per-artifact hashes and the primary checkpoint hash:
\begin{center}\small\texttt{\PrimarySourceHash}\end{center}
Aggregate data and generated tables accompany the draft. The local raw evidence
is not claimed to be a public, archived dataset, and no artifact DOI is claimed.
Before submission, the authors must complete the bibliography, authorship,
venue formatting, artifact release, and independent verification of the
manuscript. This working draft and its analysis tooling were prepared with AI
assistance and require human scientific review.

\bibliographystyle{plain}
\bibliography{references}

\appendix
\section{Model-Stratified Primary Outcomes}
\label{app:models}
Table~\ref{tab:models} preserves the model identifier distinction rather than
treating model outputs as interchangeable replications.

\begin{table}[htbp]
\centering\small
\begin{tabular}{lllrrr}
\toprule
Target & Model & Arm & $n$ & Compilation/schema pass & Visible pass \\
\midrule
\ModelResultRows
\bottomrule
\end{tabular}
\caption{Original completed results by target, model, and condition.
M1: \texttt{\ModelOne{}}; M2: \texttt{\ModelTwo{}}.}
\label{tab:models}
\end{table}

\section{Common-Intent Sensitivity}
Table~\ref{tab:matched} restricts the primary observations to intents supported
by all three targets (NET-11--NET-27). This removes one source of scenario-set
imbalance but not policy-suite differences, implementation restrictions, or
within-intent dependence.

\begin{table}[htbp]
\centering\small
\begin{tabular}{llrrr}
\toprule
Target & Arm & $n$ & Compilation/schema pass & Visible pass \\
\midrule
\MatchedResultRows
\bottomrule
\end{tabular}
\caption{Primary results restricted to the \MatchedIntentCount{} common intents,
pooled across two models. These are not new observations or independent cloud replications.}
\label{tab:matched}
\end{table}
\end{document}
